% common_setup.tex — shared objects for Part I/II formalization
% Scope: frozen writeup only. Do not import CREDO/CREME problem objects.
\documentclass[11pt]{article}
\usepackage{amsmath,amssymb,amsthm}
\usepackage[margin=1in]{geometry}
\newtheorem{assumption}{Assumption}
\newtheorem{definition}{Definition}
\newtheorem{remark}{Remark}

\title{Common setup --- Wenbin UQ (frozen)}
\author{research/formal (working)}
\date{27 July 2026}

\begin{document}
\maketitle

\section{Objects}
Let $\mathcal{D}=\{(X_i,Y_i)\}_{i=1}^n$ denote calibration data.
Let $\theta$ be a hyperparameter of an uncertainty set
$\mathcal{C}(X;\mathcal{D},\theta)$.
Miscoverage level $\alpha\in(0,1)$.

\begin{assumption}[Fixed-$\theta$ validity]
For any given $\mathcal{D}$ and \emph{prespecified} $\theta$,
\begin{equation}
\mathbb{P}\bigl(Y\in\mathcal{C}(X;\mathcal{D},\theta)\bigr)\ge 1-\alpha.
\tag{1}
\end{equation}
\end{assumption}

\begin{assumption}[Data-dependent selection program]
There exist $\hat f,\hat g$ with $\sigma(\hat f),\sigma(\hat g)\subseteq\sigma(\mathcal{D})$ and
\begin{equation}
\hat\theta(\mathcal{D})
\in
\operatorname*{arg\,min}_{\theta}
\bigl\{\hat f(y;\theta):\hat g(\theta)\le 0\bigr\}.
\tag{2}
\end{equation}
\end{assumption}

\begin{remark}[Notation quarantine]
Writeup (1) uses third argument $\theta$; prose may say $\alpha$.
Writeup uses lowercase $y$ in $\hat f(y;\theta)$.
See \texttt{research/notes/notation\_quarantine.md}.
Local convention here: treat $y$ in (2) as a placeholder argument of $\hat f$
(possibly a fixed template outcome / empirical functional of $\mathcal{D}$),
not necessarily equal to the test $Y$ in (1). Flag if lemmas need a sharper reading.
\end{remark}

\begin{remark}[Alias $C^{(q)}$]
Later Part~II files may write $C^{(q)}(X;D,\theta)$ for the writeup set
$\mathcal{C}(X;\mathcal{D},\theta)$ run at classical miscoverage level $q$
(instead of $\alpha$). This does \textbf{not} put the level into the third slot;
$\theta$ remains the optimized hyperparameter.
\end{remark}

\begin{definition}[Leave-one-out neighbor]
$\mathcal{D}^{(i)}$ is obtained from $\mathcal{D}$ by replacing the $i$-th pair
by an independent draw from the same distribution (or deleting it --- state which
when used).
\end{definition}

\begin{definition}[Selected object]
The live selected object is only $\hat\theta(\mathcal{D})$ from (2).
\end{definition}

\section{Goals (reference)}
\begin{enumerate}
  \item[(i)] When does (1) hold at $\theta=\hat\theta(\mathcal{D})$?
  \item[(ii)] If not, what is the nominal-vs-actual miscoverage deviation?
  \item[(iii)] Restore (1) at the selected hyperparameter \emph{without} recalibrating
  $\mathcal{C}$, by data-randomizing the selection map (W5).
\end{enumerate}

\end{document}
